\documentclass[a4paper,12pt]{article}

% ==================================================
% PACOTES BÁSICOS
% ==================================================

\usepackage[brazil]{babel}
\usepackage[utf8]{inputenc}
\usepackage[T1]{fontenc}

\usepackage{geometry}
\usepackage{setspace}
\usepackage{titlesec}
\usepackage{hyperref}
\usepackage{graphicx}
\usepackage{caption}
\usepackage{subcaption}
\usepackage{fancyhdr}
\usepackage{xcolor}

% Matemática
\usepackage{amsmath}
\usepackage{amssymb}

% Tabelas
\usepackage{array}
\usepackage{booktabs}

% Bibliografia
\usepackage[
    backend=biber,
    style=abnt
]{biblatex}

% Comandos personalizados
%%%%%%%%%%%%%%%%%%%%%%%%%%%%%%%%%%%%%%%%%%%%%%%%%%
% These are some new commands that may be useful 
% for paper writing in general. If other new commands
% are needed for your specific paper, please feel 
% free to add here. 
%
% The currently available commands are organized in: 
% 1) Systems
% 2) Quantities
% 3) Energies and units
% 4) particle species
% 5) Colors package
% 6) hyperlink
%%%%%%%%%%%%%%%%%%%%%%%%%%%%%%%%%%%%%%%%%%%%%%%%%%

\usepackage{amsmath}
\usepackage{amssymb}
\usepackage{upgreek}
\usepackage{multirow}
\usepackage{setspace}% http://ctan.org/pkg/setspace
\usepackage{fancyhdr}
\usepackage{datetime}

% 1) SYSTEMS
\newcommand{\btc}               {\textbf{BTC}}
\newcommand{\btcspace}          {\textbf{BTC} }
\newcommand{\pow}               {\textbf{PoW}}

\usepackage{hyperref}
\hypersetup{
    colorlinks=true, %set "true" if you want colored links
    linktoc=all,     %set to "all" if you want both sections and subsections linked
    linkcolor=blue,  %choose some color if you want links to stand out
    citecolor= blue, % color of \cite{} in the text.
    urlcolor  = blue, % color of the link for the paper in references.
}

% 5) Tikz and figures
\usepackage{epsfig}
\usepackage{lmodern}
\usepackage{mathtools}
\usepackage[utf8]{luainputenc}
\usepackage{xspace}
\usepackage{tikz}
\usepackage{pgfplots}
\pgfplotsset{compat=newest}

\usetikzlibrary{positioning}
\usepackage{subcaption}

% 6) colors:
\usepackage{xcolor}
\definecolor{ao(english)}{rgb}{0.0, 0.5, 0.0} % dark green

% 7) Add lines numbers
%\usepackage{lineno}

% add pdf file to thesis:
\usepackage{pdfpages}

\usepackage[most]{tcolorbox}

%textmarker style from colorbox doc
\tcbset{textmarker/.style={%
        enhanced,
        parbox=false,boxrule=0mm,boxsep=0mm,arc=0mm,
        outer arc=0mm,left=6mm,right=3mm,top=7pt,bottom=7pt,
        toptitle=1mm,bottomtitle=1mm,oversize}}


% define new colorboxes
\newtcolorbox{hintBox}{textmarker,
    borderline west={6pt}{0pt}{yellow},
    colback=yellow!10!white}
\newtcolorbox{importantBox}{textmarker,
    borderline west={6pt}{0pt}{red},
    colback=red!10!white}
\newtcolorbox{noteBox}{textmarker,
    borderline west={6pt}{0pt}{green},
    colback=green!10!white}

% define commands for easy access
\newcommand{\note}[1]{\begin{noteBox} \textbf{} #1 \end{noteBox}}
\newcommand{\warning}[1]{\begin{hintBox} \textbf{Warning:} #1 \end{hintBox}}
\newcommand{\important}[1]{\begin{importantBox} \textbf{} #1 \end{importantBox}}

% ==================================================
% CORES
% ==================================================

\definecolor{azul}{RGB}{30,90,150}
\definecolor{cinza}{RGB}{245,245,245}
\definecolor{verde}{RGB}{230,245,235}
\definecolor{amarelo}{RGB}{255,248,220}

% ==================================================
% CAIXAS
% ==================================================

% --------------------------------------------------
% IDEIA-CHAVE
% --------------------------------------------------

\newtcolorbox{ideia}{
    colback=azul!5,
    colframe=azul,
    title=\textbf{IDEIA-CHAVE},
    fonttitle=\bfseries,
    breakable
}

% --------------------------------------------------
% EXEMPLO
% --------------------------------------------------

\newtcolorbox{exemplo}{
    colback=verde,
    colframe=green!50!black,
    title=\textbf{EXEMPLO},
    fonttitle=\bfseries,
    breakable
}

% --------------------------------------------------
% ATENÇÃO
% --------------------------------------------------

\newtcolorbox{atencao}{
    colback=amarelo,
    colframe=orange!80!black,
    title=\textbf{ATENÇÃO},
    fonttitle=\bfseries,
    breakable
}

% ==================================================
% CABEÇALHO
% ==================================================

\pagestyle{fancy}

\fancyhf{}

\renewcommand{\headrulewidth}{0pt}

\fancyhead[R]{\thepage}

% ==================================================
% HYPERREF
% ==================================================

\hypersetup{
    colorlinks=true,
    linkcolor=blue,
    citecolor=blue,
    urlcolor=blue,
    pdftitle={Provão Paulista Seriado -- Cinemática},
    pdfauthor={André V. Silva}
}

% ==================================================
% BIBLIOGRAFIA
% ==================================================

\addbibresource{bibliography.bib}

\newcommand{\printingbibliography}{%
    \pagestyle{myheadings}
    \markright{}
    \sloppy

    \printbibliography[
        heading=bibintoc,
        title={Referências}
    ]

    \fussy
}

% ==================================================
% CONFIGURAÇÕES DE PÁGINA
% ==================================================

\renewcommand{\baselinestretch}{1.5}

\geometry{
    a4paper,
    top=30mm,
    bottom=20mm,
    left=30mm,
    right=20mm
}

% ==================================================
% FORMATAÇÃO DOS TÍTULOS
% ==================================================

\titleformat{\section}
    {\bfseries\large}
    {\thesection}
    {1em}
    {}

\titleformat{\subsection}
    {\bfseries\normalsize}
    {\thesubsection}
    {1em}
    {}

% ==================================================
% TÍTULO
% ==================================================

\title{
    \textbf{\large
    PROVÃO PAULISTA SERIADO --\\
    1ª, 2ª E 3ª SÉRIES DO ENSINO MÉDIO}
}

\author{
    André V. Silva \\[4pt]
    \href{https://andrevsilva.com/}{andrevsilva.com}
}

\date{\today}

% ==================================================
% DOCUMENTO
% ==================================================

\begin{document}

\maketitle

% ==================================================
% SUMÁRIO
% ==================================================

\tableofcontents

\newpage

% ==================================================
% CINEMÁTICA
% ==================================================

\section{Cinemática}

A Cinemática é a parte da Física que descreve os movimentos dos corpos,
analisando grandezas como posição, deslocamento, velocidade e aceleração,
sem investigar as causas responsáveis pelo movimento.

\begin{ideia}

Para estudar um movimento, é importante identificar:

\begin{itemize}

    \item a posição do corpo;

    \item o deslocamento realizado;

    \item a velocidade;

    \item a aceleração;

    \item o intervalo de tempo considerado.

\end{itemize}

\end{ideia}

% ==================================================
\subsection{Movimento Uniforme}

No \textbf{movimento uniforme (MU)}, o móvel apresenta velocidade
constante e, portanto, sua aceleração é nula:

\[
\boxed{a=0}
\]

Como a velocidade permanece constante, o móvel percorre deslocamentos
iguais em intervalos de tempo iguais.

A função horária da posição é:

\[
\boxed{x=x_0+vt}
\]

em que:

\begin{itemize}

    \item \(x\) é a posição final;

    \item \(x_0\) é a posição inicial;

    \item \(v\) é a velocidade constante;

    \item \(t\) é o tempo.

\end{itemize}

\begin{exemplo}

Um automóvel está na posição \(20\,\mathrm{m}\) e movimenta-se com
velocidade constante de \(10\,\mathrm{m/s}\).

Sua posição após \(5\,\mathrm{s}\) será:

\[
x=x_0+vt
\]

\[
x=20+10(5)
\]

\[
\boxed{x=70\,\mathrm{m}}
\]

\end{exemplo}

% ==================================================
\subsection{Movimento Uniformemente Variado}

No \textbf{movimento uniformemente variado (MUV)}, a aceleração é
constante e diferente de zero.

Isso significa que a velocidade varia de maneira uniforme ao longo
do tempo.

A função horária da velocidade é:

\[
\boxed{v=v_0+at}
\]

em que:

\begin{itemize}

    \item \(v\) é a velocidade no instante \(t\);

    \item \(v_0\) é a velocidade inicial;

    \item \(a\) é a aceleração constante;

    \item \(t\) é o tempo.

\end{itemize}

A função horária da posição é:

\[
\boxed{x=x_0+v_0t+\frac{at^2}{2}}
\]

Essa equação permite determinar a posição do móvel em determinado
instante.

\begin{exemplo}

Um móvel parte do repouso com aceleração constante de
\(2\,\mathrm{m/s^2}\).

Considerando \(x_0=0\), determine sua posição após \(4\,\mathrm{s}\).

\[
x=x_0+v_0t+\frac{at^2}{2}
\]

Como \(x_0=0\) e \(v_0=0\):

\[
x=\frac{2(4)^2}{2}
\]

\[
\boxed{x=16\,\mathrm{m}}
\]

\end{exemplo}

% ==================================================
\subsection{Equação de Torricelli}

Quando o tempo não é conhecido ou não é necessário determiná-lo,
podemos utilizar a equação de Torricelli:

\[
\boxed{v^2=v_0^2+2a\Delta x}
\]

em que:

\[
\Delta x=x-x_0
\]

Essa equação relaciona diretamente velocidade, aceleração e
deslocamento.

\begin{atencao}

A equação de Torricelli é especialmente útil em problemas de MUV
nos quais o tempo não é fornecido.

\end{atencao}

% ==================================================
\subsection{Queda Livre}

A \textbf{queda livre} é um caso particular do movimento uniformemente
variado no qual o corpo se movimenta sob a ação da gravidade,
desprezando-se a resistência do ar.

Próximo à superfície da Terra, podemos considerar:

\[
\boxed{g\approx9{,}8\,\mathrm{m/s^2}}
\]

Em muitos exercícios escolares, utiliza-se:

\[
\boxed{g\approx10\,\mathrm{m/s^2}}
\]

Quando o corpo é abandonado do repouso, temos:

\[
v_0=0
\]

Escolhendo o sentido para baixo como positivo:

\[
a=g
\]

Assim:

\[
\boxed{v=gt}
\]

e:

\[
\boxed{\Delta y=\frac{gt^2}{2}}
\]

A relação entre velocidade e deslocamento pode ser obtida pela
equação de Torricelli:

\[
\boxed{v^2=2g\Delta y}
\]

\begin{exemplo}

Uma pedra é abandonada do repouso de uma altura de \(20\,\mathrm{m}\).

Considere \(g=10\,\mathrm{m/s^2}\) e despreze a resistência do ar.

O tempo de queda pode ser determinado por:

\[
\Delta y=\frac{gt^2}{2}
\]

\[
20=\frac{10t^2}{2}
\]

\[
20=5t^2
\]

\[
t^2=4
\]

\[
\boxed{t=2\,\mathrm{s}}
\]

\end{exemplo}

% ==================================================
\subsection{Lançamento Vertical}

No \textbf{lançamento vertical}, o corpo recebe uma velocidade inicial
para cima ou para baixo e passa a sofrer a ação da gravidade.

Considerando o sentido para cima como positivo:

\[
\boxed{a=-g}
\]

A função horária da velocidade é:

\[
\boxed{v=v_0-gt}
\]

A função horária da posição é:

\[
\boxed{
y=y_0+v_0t-\frac{gt^2}{2}
}
\]

A equação de Torricelli fica:

\[
\boxed{
v^2=v_0^2-2g\Delta y
}
\]

No ponto de altura máxima, a velocidade instantânea é nula:

\[
\boxed{v=0}
\]

Portanto, podemos determinar o tempo de subida por:

\[
0=v_0-gt
\]

\[
\boxed{
t_{\mathrm{subida}}=\frac{v_0}{g}
}
\]

A altura máxima, em relação ao ponto de lançamento, pode ser obtida
por:

\[
\boxed{
h_{\mathrm{máx}}=\frac{v_0^2}{2g}
}
\]

\begin{exemplo}

Uma bola é lançada verticalmente para cima com velocidade inicial
de \(20\,\mathrm{m/s}\).

Considere \(g=10\,\mathrm{m/s^2}\).

O tempo necessário para atingir a altura máxima é:

\[
t_{\mathrm{subida}}=\frac{v_0}{g}
\]

\[
t_{\mathrm{subida}}=\frac{20}{10}
\]

\[
\boxed{
t_{\mathrm{subida}}=2\,\mathrm{s}
}
\]

A altura máxima atingida será:

\[
h_{\mathrm{máx}}
=
\frac{20^2}{2(10)}
\]

\[
\boxed{
h_{\mathrm{máx}}=20\,\mathrm{m}
}
\]

\end{exemplo}

% ==================================================
\subsection{Relação entre as Funções do Movimento}

As principais funções utilizadas na Cinemática podem ser organizadas
da seguinte maneira:

\begin{center}

\renewcommand{\arraystretch}{1.6}

\begin{tabular}{|>{\centering\arraybackslash}m{4cm}|
                >{\centering\arraybackslash}m{8cm}|}

\hline

\textbf{Movimento} &
\textbf{Equação} \\

\hline

Movimento Uniforme &
\(x=x_0+vt\) \\

\hline

MUV -- velocidade &
\(v=v_0+at\) \\

\hline

MUV -- posição &
\(x=x_0+v_0t+\frac{at^2}{2}\) \\

\hline

Equação de Torricelli &
\(v^2=v_0^2+2a\Delta x\) \\

\hline

Queda livre &
\(v=gt\) \\

\hline

Queda livre &
\(\Delta y=\frac{gt^2}{2}\) \\

\hline

Lançamento vertical &
\(v=v_0-gt\) \\

\hline

Lançamento vertical &
\(y=y_0+v_0t-\frac{gt^2}{2}\) \\

\hline

\end{tabular}

\end{center}

% ==================================================
\subsection{Resumo Conceitual}

\begin{ideia}

\textbf{Movimento Uniforme (MU)}

\[
a=0
\]

\[
v=\text{constante}
\]

\[
x=x_0+vt
\]

\vspace{0.3cm}

\textbf{Movimento Uniformemente Variado (MUV)}

\[
a=\text{constante}
\]

\[
v=v_0+at
\]

\[
x=x_0+v_0t+\frac{at^2}{2}
\]

\vspace{0.3cm}

\textbf{Queda Livre}

\[
v_0=0
\]

\[
a=g
\]

\[
v=gt
\]

\[
\Delta y=\frac{gt^2}{2}
\]

\vspace{0.3cm}

\textbf{Lançamento Vertical para Cima}

\[
a=-g
\]

\[
v=v_0-gt
\]

\[
y=y_0+v_0t-\frac{gt^2}{2}
\]

No ponto mais alto:

\[
v=0
\]

\end{ideia}

% ==================================================
\subsection{Cuidados com os Sinais}

\begin{atencao}

O sinal da aceleração depende da orientação do eixo escolhido.

No lançamento vertical, se o sentido positivo for escolhido para cima,
a aceleração da gravidade será:

\[
a=-g
\]

Na queda livre, se o sentido positivo for escolhido para baixo:

\[
a=+g
\]

O mais importante é manter a mesma convenção de sinais durante toda
a resolução do problema.

\end{atencao}

% ==================================================
\subsection{Estratégia para Resolver Problemas}

Em questões de Cinemática, uma estratégia eficiente é:

\begin{enumerate}

    \item Identificar o tipo de movimento.

    \item Listar os dados fornecidos pelo problema.

    \item Escolher o eixo e definir a convenção de sinais.

    \item Identificar a grandeza que se deseja determinar.

    \item Escolher a equação adequada.

    \item Substituir os valores, mantendo as unidades consistentes.

    \item Verificar se o resultado possui unidade e sentido físico
    coerentes.

\end{enumerate}

\begin{atencao}

Antes de utilizar uma equação, verifique se ela corresponde ao tipo
de movimento estudado.

Por exemplo:

\begin{itemize}

    \item velocidade constante $\rightarrow$ Movimento Uniforme;

    \item aceleração constante $\rightarrow$ Movimento Uniformemente
    Variado;

    \item movimento vertical sob ação da gravidade $\rightarrow$
    utilizar \(g\) como aceleração.

\end{itemize}

\end{atencao}

% ==================================================
% BIBLIOGRAFIA
% ==================================================

\printingbibliography

% ==================================================
% FIM DO DOCUMENTO
% ==================================================

\end{document}